\ifdefined\gkver\else\def\gkver{student}\fi
\documentclass[\gkver]{gaokaozhenti}
\begin{document}

\chapter{2014年浙江理综}
%% paperType: 理综物理部分

\begin{enumerate}
\item
%% number: 14
%% paperName: 2014年浙江理综第 14 题 4 分
%% typeId: 1
%% type: 单选题
%% chapter: 机械波
%% point: 波长频率波速
%% method: 无
%% score: 4
%% degree: 650
%% duplicateId: 0
%% body:
下列说法正确的是\xzanswer{B}

\fourchoices[answer=B]
{机械波的振幅与波源无关}
{机械波的传播速度由介质本身的性质决定}
{物体受到的静摩擦力方向与其运动方向相反}
{动摩擦因数的数值跟相互接触的两个物体的材料无关}

%% answer: B
%% memo:
\memoanswer{
机械波振幅与波源有关，波速由介质决定，静摩擦力方向与相对运动趋势方向相反，与物体运动方向无关，动摩擦因数与接触物材料、粗糙程度、接触面温度等有关，所以选 B。
}

\item
%% number: 15
%% paperName: 2014年浙江理综第 15 题 4 分
%% typeId: 1
%% type: 单选题
%% chapter: 交流电
%% point: 变压器
%% method: 无
%% score: 4
%% degree: 650
%% duplicateId: 0
%% body:
如图所示为远距离交流输电的简化电路图。发电厂的输出电压是 $U$，用等效总电阻是 $r$ 的两条输电线输电，输电线路中的电流是 $I_{1}$，其末端间的电压为 $U_{1}$。在输电线与用户间连有一个理想变压器，流入用户端的电流是 $I_{2}$。则\xzanswer{A}

\onepicture[w=9cm, align=right]{figs/15.jpg}

\fourchoices[answer=A]
{用户端的电压为 $I_{1}U_{1}/I_{2}$}
{输电线上的电压降为 $U$}
{理想变压器的输入功率为 $I_{1}^{2}r$}
{输电线路上损失的电功率为 $I_{1}U$}

%% answer: A
%% memo:
\memoanswer{
理想变压器输入端与输出端功率相等，$U_{1}I_{1}=U_{2}I_{2}$，用户端的电压 $U_{2}=\frac{I_{1}}{I_{2}}U_{1}$，选项 A 正确；输电线上的电压降 $\Delta U=U-U_{1}=I_{1}r$，选项 B 错误；理想变压器输电线上损失的功率为 $I_{1}^{2}r$，选项 C、D 错误。
}

\item
%% number: 16
%% paperName: 2014年浙江理综第 16 题 4 分
%% typeId: 1
%% type: 单选题
%% chapter: 曲线运动
%% point: 开普勒定律
%% method: 无
%% score: 4
%% degree: 450
%% duplicateId: 0
%% body:
长期以来“卡戎星（Charon）”被认为是冥王星唯一的卫星，它的公转轨道半径 $r_{1}=19600\Ukm$，公转周期 $T_{1}=6.39$ 天。2006 年 3 月，天文学家新发现两颗冥王星的小卫星，其中一颗的公转轨道半径 $r_{2}=48000\Ukm$，则它的公转周期 $T_{2}$ 最接近于\xzanswer{B}

\fourchoices[answer=B]
{15 天}
{25 天}
{35 天}
{45 天}

%% answer: B
%% memo:
\memoanswer{
根据开普勒第三定律 $\frac{r_{1}^{3}}{T_{1}^{2}}=\frac{r_{2}^{3}}{T_{2}^{2}}$，代入数据计算可得 $T_{2}$ 约等于 25 天。选项 B 正确。
}

\item
%% number: 17
%% paperName: 2014年浙江理综第 17 题 4 分
%% typeId: 1
%% type: 单选题
%% chapter: 机械振动
%% point: 振动图象
%% method: 无
%% score: 4
%% degree: 450
%% duplicateId: 0
%% body:
一位游客在千岛湖边欲乘游船，当日风浪很大，游船上下浮动。可把游船浮动简化成竖直方向的简谐运动，振幅为 $20\Ucm$，周期为 $3.0\Us$。当船上升到最高点时，甲板刚好与码头地面平齐。地面与甲板的高度差不超过 $10\Ucm$ 时，游客能舒服地登船。在一个周期内，游客能舒服地登船的时间是\xzanswer{C}

\fourchoices[answer=C]
{$0.5\Us$}
{$0.75\Us$}
{$1.0\Us$}
{$1.5\Us$}

%% answer: C
%% memo:
\memoanswer{
从平衡位置开始计时，游船的振动方程 $x=20\sin\frac{2\pi}{3}t\Ucm$，游客要舒服地登船需满足的条件 $\Delta x=20-x\leqslant10$，解得 $0.25\Us\leqslant t\leqslant1.25\Us$，故游客能舒服地登船的时间 $\Delta t=1.0\Us$，选项 C 正确。
}

\item
%% number: 18
%% paperName: 2014年浙江理综第 18 题 6 分
%% typeId: 2
%% type: 多选题
%% chapter: 光学
%% point: 光的干涉
%% method: 无
%% score: 6
%% degree: 450
%% duplicateId: 0
%% body:
关于下列光学现象，说法正确的是\xzanswer{CD}

\fourchoices[answer=CD]
{水中蓝光的传播速度比红光快}
{光从空气射入玻璃时可能发生全反射}
{在岸边观察前方水中的一条鱼，鱼的实际深度比看到的要深}
{分别用蓝光和红光在同一装置上做双缝干涉实验，用红光时得到的条纹间距更宽}

%% answer: CD
%% memo:
\memoanswer{
在同一种介质中，波长越短，波速越慢，故红光的传播速度比蓝光大，选项 A 错误；光从空气射向玻璃是从光疏介质射向光密介质，不能发生全反射，选项 B 错误；在岸边观察水中的鱼，由于光的折射，鱼的实际深度比看到的深度要深，选项 C 正确；在空气中红光的波长比蓝光要长，根据 $\Delta x=\frac{l}{d}\lambda$ 可知红光的双缝干涉条纹间距大，选项 D 正确。
}

\item
%% number: 19
%% paperName: 2014年浙江理综第 19 题 6 分
%% typeId: 2
%% type: 多选题
%% chapter: 电场
%% point: 带电粒子的平衡
%% method: 无
%% score: 6
%% degree: 650
%% duplicateId: 0
%% body:
如图所示，水平地面上固定一个光滑绝缘斜面，斜面与水平面的夹角为 $\theta$。一根轻质绝缘细线的一端固定在斜面顶端，另一端系有一个带电小球 A，细线与斜面平行。小球 A 的质量为 $m$、电量为 $q$。小球 A 的右侧固定放置带等量同种电荷的小球 B，两球心的高度相同、间距为 $d$。静电力常量为 $k$，重力加速度为 $g$，两带电小球可视为点电荷。小球 A 静止在斜面上，则\xzanswer{AC}

\onepicture[w=9cm, align=right]{figs/19.png}

\fourchoices[answer=AC]
{小球 A 与 B 之间库仑力的大小为 $\frac{kq^{2}}{d^{2}}$}
{当 $\frac{q}{d}=\sqrt{\frac{mg\sin\theta}{k}}$ 时，细线上的拉力为 0}
{当 $\frac{q}{d}=\sqrt{\frac{mg\tan\theta}{k}}$ 时，细线上的拉力为 0}
{当 $\frac{q}{d}=\sqrt{\frac{mg}{k\tan\theta}}$ 时，斜面对小球 A 的支持力为 0}

%% answer: AC
%% memo:
\memoanswer{
点电荷库仑定律 $F=\frac{kq^{2}}{d^{2}}$，所以 A 正确；当细线上的拉力为 0 的时候，小球 A 受到库仑力、斜面支持力、重力，具体关系为 $\frac{kq^{2}}{d^{2}}=mg\tan\theta$，即 C 选项正确。由受力分析可知，斜面对小球的支持力不可能为 0，所以 D 选项错误。
}

\item
%% number: 20
%% paperName: 2014年浙江理综第 20 题 6 分
%% typeId: 2
%% type: 多选题
%% chapter: 电磁感应
%% point: 电磁感应受力分析
%% method: 无
%% score: 6
%% degree: 450
%% duplicateId: 0
%% body:
如图 1 所示，两根光滑平行导轨水平放置，间距为 $L$，其间有竖直向下的匀强磁场，磁感应强度为 $B$。垂直于导轨水平对称放置一根均匀金属棒。从 $t=0$ 时刻起，棒上有如图 2 所示的持续交流电流 $I$，周期为 $T$，最大值为 $I_{m}$，图 1 中 $I$ 所示方向为电流正方向。则金属棒\xzanswer{ABC}

\onepicture[w=11cm, align=right]{figs/20.png}

\fourchoices[answer=ABC]
{一直向右移动}
{速度随时间周期性变化}
{受到的安培力随时间周期性变化}
{受到的安培力在一个周期内做正功}

%% answer: ABC
%% memo:
\memoanswer{
根据题意得出 $v-t$ 图象如图所示，金属棒一直向右运动，A 正确。速度随时间做周期性变化，B 正确。据 $F_{\text{安}}=BIL$ 及左手定则可判定，$F_{\text{安}}$ 大小不变，方向做周期性变化，则 C 项正确。$F_{\text{安}}$ 在前半周期做正功，后半周期做负功，则 D 项错。

\onepicture[w=6cm]{figs/20_an.jpg}
}

\item
%% number: 21
%% paperName: 2014年浙江理综第 21 题 10 分
%% typeId: 4
%% type: 实验题
%% chapter: 力物体的平衡
%% point: 胡克定律
%% method: 无
%% score: 10
%% degree: 650
%% duplicateId: 0
%% body:
在“探究弹力和弹簧伸长的关系”时，某同学把两根弹簧如图 1 连接起来进行探究。

\begin{center}
\begin{tabular}{|c|c|c|c|c|}
\hline
钩码数 & 1 & 2 & 3 & 4 \\ \hline
$L_{\mathrm{A}}/\Ucm$ & 15.71 & 19.71 & 23.66 & 27.76 \\ \hline
$L_{\mathrm{B}}/\Ucm$ & 29.96 & 35.76 & 41.51 & 47.36 \\ \hline
\end{tabular}
\end{center}

\begin{enumerate}
	\item
	某次测量如图 2 所示，指针示数为\_\_\_\_\_\_\_\_\_\_\_$\Ucm$。

	\item
	在弹性限度内，将 $50\Ug$ 的钩码逐个挂在弹簧下端，得到指针 A、B 的示数 $L_{\mathrm{A}}$ 和 $L_{\mathrm{B}}$ 如表 1。用表 1 数据计算弹簧 Ⅰ 的劲度系数为\_\_\_\_\_\_\_\_\_\_\_$\UNm$（重力加速度 $g=10\Umsq$）。由表 1 数据\_\_\_\_\_（填“能”或“不能”）计算出弹簧 Ⅱ 的劲度系数。
\end{enumerate}

\twopicture[showlabel=false, align=right]
{figs/21b.jpg}
{figs/21a.jpg}

%% answer:
\jdanswer{
\begin{enumerate}
	\item
	$15.95\Ucm\sim16.05\Ucm$，有效数字位数正确

	\item
	$12.2\UNm\sim12.8\UNm$，能

\end{enumerate}
}
%% memo:
\memoanswer{
（1）由图 2 可知刻度尺能精确到 $0.1\Ucm$，读数时需要往后估读一位。故指针示数为 $16.00\pm0.05\Ucm$。

（2）由表 1 中数据可知每挂一个钩码，弹簧 Ⅰ 的平均伸长量 $\Delta x_{1}\approx4\Ucm$，弹簧 Ⅱ 的总平均伸长量 $\Delta x_{2}\approx5.80\Ucm$，根据胡克定律可求得弹簧 Ⅰ 的劲度系数为 $12.5\UNm$，同理也能求出弹簧 Ⅱ 的劲度系数。
}

\item
%% number: 22
%% paperName: 2014年浙江理综第 22 题 10 分
%% typeId: 4
%% type: 实验题
%% chapter: 电路
%% point: 伏安法测电阻
%% method: 无
%% score: 10
%% degree: 650
%% duplicateId: 0
%% body:
小明对 2B 铅笔芯的导电性能感兴趣，于是用伏安法测量其电阻值。

\begin{enumerate}
	\item
	图 1 是部分连接好的实物电路图，请用电流表外接法完成接线并在图 1 中画出。

	\item
	小明用电流表内接法和外接法分别测量了一段 2B 铅笔芯的伏安特性，并将得到的电流、电压数据描到 $U-I$ 图上，如图 2 所示。在图中，由电流表外接法得到的数据点是用\_\_\_\_\_\_\_（填“$\bigcirc$”或“$\times$”）表示的。

	\item
	请你选择一组数据点，在图 2 上用作图法作图，并求出这段铅笔芯的电阻为\_\_\_\_\_\_\_$\UO$。
\end{enumerate}

\twopicture[showlabel=false, align=right]
{figs/22a.jpg}
{figs/22b.jpg}

%% answer:
\jdanswer{
\begin{enumerate}
	\item
	如图所示：

	\item
	“$\times$”

	\item
	作图如下图：

	用“$\times$”连线 $R=(1.1\sim1.3)\UO$；用“$\bigcirc$”连线 $R=(1.5\sim1.7)\UO$

\end{enumerate}

\onepicture[w=6cm]{figs/22_an1.jpg}

\onepicture[w=7cm]{figs/22_an2.jpg}
}
%% memo:
\memoanswer{
（1）由图 2 中的电压、电流数据从零开始可知滑动变阻器采用分压式，电压表选择量程 $3\UV$，电流表采用外接法。

（2）外接法由于电压表分流，测得的电阻比内接法测得的要小，故电流表外接法得到的数据点是用“$\times$”表示的。

（3）用“$\times$”数据点连直线，斜率为铅笔芯的电阻，考虑误差因素，$R=(1.1\sim1.3)\UO$，用“$\bigcirc$”数据点连直线，同理得 $R=(1.5\sim1.7)\UO$。
}

\item
%% number: 23
%% paperName: 2014年浙江理综第 23 题 16 分
%% typeId: 6
%% type: 计算题
%% chapter: 曲线运动
%% point: 平抛运动
%% method: 无
%% score: 16
%% degree: 450
%% duplicateId: 0
%% body:
如图所示，装甲车在水平地面上以速度 $v_{0}=20\Ums$ 沿直线前进，车上机枪的枪管水平，距地面高为 $h=1.8\Um$。在车正前方竖直放置一块高为两米的长方形靶，其底边与地面接触。枪口与靶距离为 $L$ 时，机枪手正对靶射出第一发子弹，子弹相对于枪口的初速度为 $v=800\Ums$。在子弹射出的同时，装甲车开始匀减速运动，行进 $s=90\Um$ 后停下。装甲车停下后，机枪手以相同方式射出第二发子弹。（不计空气阻力，子弹看成质点，重力加速度 $g=10\Umsq$）

\begin{enumerate}
	\item
	求装甲车匀减速运动时的加速度大小；

	\item
	当 $L=410\Um$ 时，求第一发子弹的弹孔离地的高度，并计算靶上两个弹孔之间的距离；

	\item
	若靶上只有一个弹孔，求 $L$ 的范围。
\end{enumerate}

\onepicture[w=12cm, align=right]{figs/23.jpg}

%% answer:
\jdanswer{
\begin{enumerate}
	\item
	$\frac{20}{9}\Umsq$

	\item
	$0.55\Um\sim0.45\Um$

	\item
	$492\Um<L\leq570\Um$

\end{enumerate}
}
%% memo:
\memoanswer{
（1）装甲车加速度 $a=\frac{v_{0}^{2}}{2s}=\frac{20}{9}\Umsq$。

（2）第一发子弹飞行时间 $t_{1}=\frac{L}{v+v_{0}}=0.5\Us$，

弹孔离地高度 $h_{1}=h-\frac{1}{2}gt_{1}^{2}=0.55\Um$，

第二发子弹离地的高度 $h_{2}=h-\frac{1}{2}g\left(\frac{L-s}{v}\right)^{2}=1.0\Um$，

两弹孔之间的距离 $\Delta h=h_{2}-h_{1}=0.45\Um$。

（3）第一发子弹打到靶的下沿时，装甲车离靶的距离为 $L_{1}$，

$L_{1}=(v_{0}+v)\sqrt{\frac{2h}{g}}=492\Um$；

第二发子弹打到靶的下沿时，装甲车离靶的距离为 $L_{2}$，

$L_{2}=v\sqrt{\frac{2h}{g}}+s=570\Um$。

$L$ 的范围 $492\Um<L\leq570\Um$。
}

\item
%% number: 24
%% paperName: 2014年浙江理综第 24 题 20 分
%% typeId: 6
%% type: 计算题
%% chapter: 电磁感应
%% point: 电磁感应受力分析
%% method: 无
%% score: 20
%% degree: 450
%% duplicateId: 0
%% body:
某同学设计一个发电测速装置，工作原理如图所示。一个半径为 $R=0.1\Um$ 的圆形金属导轨固定在竖直平面上，一根长为 $R$ 的金属棒 OA，A 端与导轨接触良好，O 端固定在圆心处的转轴上。转轴的左端有一个半径为 $r=R/3$ 的圆盘，圆盘和金属棒能随转轴一起转动。圆盘上绕有不可伸长的细线，下端挂着一个质量为 $m=0.5\Ukg$ 的铝块。在金属导轨区域内存在垂直于导轨平面向右的匀强磁场，磁感应强度 $B=0.5\UT$。a 点与导轨相连，b 点通过电刷与 O 端相连。测量 a、b 两点间的电势差 $U$ 可算得铝块速度。铝块由静止释放，下落 $h=0.3\Um$ 时，测得 $U=0.15\UV$。（细线与圆盘间没有滑动，金属棒、导轨、导线及电刷的电阻均不计，重力加速度 $g=10\Umsq$）

\begin{enumerate}
	\item
	测 $U$ 时，a 点相接的是电压表的“正极”还是“负极”？

	\item
	求此时铝块的速度大小；

	\item
	求此下落过程中铝块机械能的损失。
\end{enumerate}

\onepicture[w=10cm, align=right]{figs/24.jpg}

%% answer:
\jdanswer{
\begin{enumerate}
	\item
	正极

	\item
	$2\Ums$

	\item
	$0.5\UJ$

\end{enumerate}
}
%% memo:
\memoanswer{
（1）正极。

（2）由电磁感应定律得 $U=E=\frac{\Delta\Phi}{\Delta t}$，

$\Delta\Phi=\frac{1}{2}BR^{2}\Delta\theta$，

$U=\frac{1}{2}B\omega R^{2}$，

$v=r\omega=\frac{1}{3}\omega R$，

所以 $v=\frac{2U}{3BR}=2\Ums$。

（3）$\Delta E=mgh-\frac{1}{2}mv^{2}$，

$\Delta E=0.5\UJ$。
}

\item
%% number: 25
%% paperName: 2014年浙江理综第 25 题 22 分
%% typeId: 6
%% type: 计算题
%% chapter: 磁场
%% point: 带电粒子在磁场中的运动
%% method: 无
%% score: 22
%% degree: 350
%% duplicateId: 0
%% body:
离子推进器是太空飞行器常用的动力系统，某种推进器设计的简化原理如图 1 所示，截面半径为 $R$ 的圆柱腔分为两个工作区。Ⅰ 为电离区，将氙气电离获得 1 价正离子；Ⅱ 为加速区，长度为 $L$，两端加有电压，形成轴向的匀强电场。Ⅰ 区产生的正离子以接近 0 的初速度进入 Ⅱ 区，被加速后以速度 $v_{M}$ 从右侧喷出。Ⅰ 区内有轴向的匀强磁场，磁感应强度大小为 $B$，在离轴线 $R/2$ 处的 C 点持续射出一定速率范围的电子。假设射出的电子仅在垂直于轴线的截面上运动，截面如图 2 所示（从左向右看）。电子的初速度方向与中心 O 点和 C 点的连线成 $\alpha$ 角（$0<\alpha<90^{\circ}$）。推进器工作时，向 Ⅰ 区注入稀薄的氙气。电子使氙气电离的最小速率为 $v_{0}$，电子在 Ⅰ 区内不与器壁相碰且能到达的区域越大，电离效果越好。已知离子质量为 $M$；电子质量为 $m$，电量为 $e$。（电子碰到器壁即被吸收，不考虑电子间的碰撞）。

\begin{enumerate}
	\item
	求 Ⅱ 区的加速电压及离子的加速度大小；

	\item
	为取得好的电离效果，请判断 Ⅰ 区中的磁场方向（按图 2 说明是“垂直纸面向里”或“垂直纸面向外”）；

	\item
	$\alpha$ 为 $90^{\circ}$ 时，要取得好的电离效果，求射出的电子速率 $v$ 的范围；

	\item
	要取得好的电离效果，求射出的电子最大速率 $v_{max}$ 与 $\alpha$ 的关系。
\end{enumerate}

\onepicture[w=12cm, align=right]{\includesvg{figs/25}}

%% answer:
\jdanswer{
\begin{enumerate}
	\item
	$\frac{v_{M}^{2}}{2L}$

	\item
	垂直纸面向外

	\item
	$v_{0}\leq v\leq\frac{3eBR}{4m}$

	\item
	$v_{max}=\frac{3eBR}{4m(2-\sin\alpha)}$

\end{enumerate}
}
%% memo:
\memoanswer{
（1）由动能定理得 $\frac{1}{2}Mv_{M}^{2}=eU$①，

$U=\frac{Mv_{M}^{2}}{2e}$②，

$a=\frac{eE}{M}=e\frac{U}{ML}=\frac{v_{M}^{2}}{2L}$③。

（2）垂直纸面向外④。

（3）设电子运动的最大半径为 $r$，

$2r=\frac{3}{2}R$⑤，

$eBv=m\frac{v^{2}}{r}$⑥，

\onepicture[w=4.5cm]{figs/25_an1.png}

所以有 $v_{0}\leqslant v\leqslant\frac{3eBR}{4m}$⑦。

要使⑦式有解，磁感应强度 $B>\frac{4mv_{0}}{3eR}$⑧。

（4）如图所示，$OA=R-r$，$OC=\frac{R}{2}$，$AC=r$，

\onepicture[w=4.5cm]{figs/25_an2.png}

根据几何关系得 $r=\frac{3R}{4(2-\sin\alpha)}$⑨，

由⑥⑨式得 $v_{max}=\frac{3eBR}{4m(2-\sin\alpha)}$。
}

\end{enumerate}

\end{document}
